% -------------------------------------------------
% PAGE DE TITRE
% -------------------------------------------------

%
% ATTENTION : ne pas utiliser les macros \title \author ou \date !!!
%
\auteur{Maxime Beauchesne-Jolin}
\titre{Un instrument d'aide à la rédaction pour les documents exprimant des modèles logiciels}

% n'existe que si la thèse est rédigée en anglais (pas obligatoire)
\englishTitle{Title in english}


% n'existent que si l'option cotutelle est présente
\institutionPartenaire{Université XYZ} 
\departementPartenaire{au département (centre, institut) ABC} 
\gradePartenaire{docteur en informatique} 

%\dedicace{À quelqu'un de significant} %%% ajouter un item de dédicace

% -------------------------------------------------

% DESCRIPTION DU SOMMAIRE (EN FRANCAIS) -----------
\sommaire{
Les problèmes en lien avec la documentation sont récurrents en informatique appliquée. Cette situation soulève de nombreuses interrogations qui conduisent à revenir aux fondamentaux. L'informatique appliquée comprend la création et la modification de modèles. Ces modèles, du moins, leurs expressions, sont la plupart du temps disséminés à travers la documentation informatique. Puisque les documents contiennent les modèles, ils sont les piliers des systèmes et des logiciels développés. Cependant, il faut reconnaître que le développement de la documentation informatique ne s'est pas réalisé avec la même rigueur et le même succès que l'algorithmique ou la modélisation de données. Ce mémoire tente d'en donner les raisons et de proposer une solution.
\\

La solution décrite dans ce mémoire répond au souhait d'avoir une documentation partageant certaines des caractéristiques de qualité des logiciels. Cela est possible dès lors que la documentation comme le logiciel concrétisent des modèles. Ainsi, plusieurs approches provenant de la documentation informatique et de la modélisation informatique sont utilisées. Il y a l'utilisation, entre autres, de la programmation littéraire, de langages formels et de méthodes formelles. Ainsi, la solution obtenue prend la forme d'un instrument d'aide à la rédaction. Cet instrument permettra aux utilisateurs de maintenir la documentation cohérente (et donc le modèle) et de l'adapter aux différents changements qui peuvent survenir lors d'un procédé de développement ou après.
}

%Provenant de https://dl.acm.org/ccs
\motsCles{documentation,
méthodes formelles,
définitions de langages formels,
compilateurs,
modèles de systèmes logiciels,
langages de description de systèmes,
méthodologies de modélisation,
ingénierie logicielle pilotée par les modèles
}

%  optionnel, valide si votre document est en anglais 
%  optionnal, valid only if your document is written in english, you have to provide a french abstract 
%  option	english du document
\abstract{Write here your abstract in english}
\keywords{
documentation,
formal methods,
formal language definitions,
compilers,
software system models,
context specific languages,
system description languages,
modeling methodologies,
model-driven software engineering
}


% REMERCIEMENTS -----------------------------------
\remerciements{Ce document est en partie le résultat d'un long processus dans lequel j'ai eu de nombreuses réflexions sur l'informatique. Si mes nombreuses années d'expérience en informatique n'ont pas entamé ma curiosité, elles ont en revanche attisé mes suspicions envers le domaine. Ainsi, j'espère que ce travail permettra au moins de ne plus tergiverser sur l'importance que peut avoir la documentation pour l'informatique appliquée.
\\

Lors de ces différentes réflexions, j'ai eu l'immense privilège d'être accompagné par le professeur Luc Lavoie. À vrai dire, il est rare de pouvoir côtoyer un érudit et de profiter de son savoir, aussi je ne saurai pas de quelle manière lui témoigner toute ma reconnaissance. Je désire aussi souligner l'effort de Maryse Couture à la révision linguistique de ce document, cela n'a probablement pas été de tout repos.
\\

Je souhaite également exprimer ma gratitude à tous les professionnels avec lesquels j'ai pu échanger sur l'informatique, ainsi qu'aux  professeurs et chargés de cours qui ont généreusement partagé avec moi leurs connaissances et leur lucidité. Vous m'avez tous permis de nourrir mes réflexions.
\\

Je ne pourrais conclure ces remerciements sans souligner l'incommensurable patience de toute ma famille (Isabelle, Zachary, Anthony et Mélodie) qui m'a permis paisiblement de m'instruire et de rédiger ce mémoire. 
}

% LISTE DES ABREVIATIONS --------------------------
%% Insérez ici  toutes les abréviations que vous souhaitez utiliser dans le document.
% les entrées seront triées

%\newacronym{etiquette}{ABR}{Expression à remplacer}
%\newacronym{edp}{EDP}{Équations aux dérivees partielles}
%\newacronym{sw}{SW}{Star Wars}
%\newacronym{st}{ST}{Star Trek}
%\newacronym{bv}{BV}{Babylon V}

%nom commun
\newabbreviation
 [type=\acronymtype]
 {dfd}{DFD}{Diagramme de flux de données}

\newabbreviation
[type=\acronymtype]
{iar}{IAR}{Instrument d'aide à la rédaction}

%nom commun
\newabbreviation
[type=\acronymtype]
{sel}{SEL}{Spécification des exigences logicielles}

%nom commun
%\newabbreviation
%[type=\acronymtype]
%{sgbd}{SGBD}{Système de gestion de base de données}

%nom commun
\newabbreviation
[type=\acronymtype]
{ae}{EA}{Entité-Association (\textit{Entity-Relationship})}

%\newabbreviation
%[type=\acronymtype]
%{cnp}{CNP}{Classification nationale des professions}




\newabbreviation
[type=\acronymtype]
{antlr}{ANTLR}{Un autre outil pour la reconnaissance du langage (\textit{ANother Tool for Language Recognition})}

\newabbreviation
[type=\acronymtype]
{bpmn}{BPMN}{Modèle de procédé d'affaire et notation (\textit{Business Process Model and Notation})}

\newabbreviation
[type=\acronymtype]
{cmmi}{CMMI}{Modèle intégré de maturité des capacités d'intégration (\textit{Capability Maturity Model Integration})}

%\newabbreviation
%[type=\acronymtype]
%{DID}{DID}{\textit{Data Item Description}}

\newabbreviation
[type=\acronymtype]
{ebnf}{EBNF}{Forme étendue de Backus-Naur (\textit{Extended Backus–Naur form})}

%nom commun
\newabbreviation
[type=\acronymtype]
{gml}{GML}{Grand modèle de langage}

%\newabbreviation
%[type=\acronymtype]
%{loc}{LOC}{\textit{Line of Code}}

%\newabbreviation
%[type=\acronymtype]
%{nasa}{NASA}{\textit{National Aeronautics and Space Administration}}

%\newabbreviation
%[type=\acronymtype]
%{sei}{SEI}{\textit{Softtware Engineering Institute}}

%\newabbreviation
%[type=\acronymtype]
%{usaf}{USAF}{\textit{United States Air Force}}

%\newabbreviation
%[type=\acronymtype]
%{mcas}{MCAS}{\textit{Maneuvering Characteristics Augmentation System}}

\newabbreviation
[type=\acronymtype]
{pmbok}{PMBOK}{Référentiel des connaissances en gestion de projet (\textit{Project Management Body of Knowledge})}

%\newabbreviation
%[type=\acronymtype]
%{itil}{ITIL}{\textit{Information Technology Infrastructure Library}}

\newabbreviation
[type=\acronymtype]
{xml}{XML}{Langage de balisage extensible (\textit{Extensible Markup Language})}

\newabbreviation
[type=\acronymtype]
{xslt}{XSLT}{Langage de transformation XML (\textit{eXtensible Stylesheet Language Transformations})}

\newabbreviation
[type=\acronymtype]
{dtd}{DTD}{Définition de type de document (\textit{Document Type Definitions})}

%\newabbreviation
%[type=\acronymtype]
%{sgml}{SGML}{\textit{Standard Generalized Markup Language}}

%\newabbreviation
%[type=\acronymtype]
%{dsdl}{DSDL}{\textit{Document Schema Definition Languages}}

%\newabbreviation
%[type=\acronymtype]
%{owl}{OWL}{\textit{Web Ontology Language}}

\newabbreviation
[type=\acronymtype]
{xsd}{XSD}{Définition du schéma XML (\textit{XML Schema Definition})}

\newabbreviation
[type=\acronymtype]
{dita}{DITA}{Architecture des types d'informations Darwin (\textit{Darwin Information Type Architecture})}

%nom commun
\newabbreviation
[type=\acronymtype]
{aa}{AA}{Apprentissage automatique}

%nom commun
\newabbreviation
[type=\acronymtype]
{rn}{RN}{Réseau de neurones}

%\newabbreviation
%[type=\acronymtype]
%{ai}{AI}{\textit{Artificial Intelligence}}

%nom commun
\newabbreviation
[type=\acronymtype]
{tal}{TAL}{Traitement automatique des langues}

%nom commun
\newabbreviation
[type=\acronymtype]
{ap}{AP}{Apprentissage profond}

%\newabbreviation
%[type=\acronymtype]
%{acl}{ACL}{\textit{Association for Computational Linguistics}}

\newabbreviation
[type=\acronymtype]
{swebok}{SWEBOK}{Guide du corpus de connaissances en génie logiciel (\textit{Software Engineering Body of Knowledge})}

\newabbreviation
[type=\acronymtype]
{uml}{UML}{Langage de modélisation unifié (\textit{Unified Modeling Language})}

\newabbreviation
[type=\acronymtype]
{sysml}{SysML}{Langage de modélisation des systèmes (\textit{SYStems Modeling Language})}

%nom commun
\newabbreviation
[type=\acronymtype]
{dsl}{DSL}{Langage spécifique à un domaine (\textit{Domain-Specific Language})}

\newabbreviation
[type=\acronymtype]
{mbebok}{MBEBOK}{Corpus de connaissances en génie logiciel basé sur les modèles (\textit{Model-Based Software Engineering Body of Knowledge})}

%nom commun
\newabbreviation
[type=\acronymtype]
{casetool}{CASE tool}{Outil de génie logiciel assisté par ordinateur (\textit{Computer-Aided Software Engineering tool})}

\newabbreviation
[type=\acronymtype]
{mda}{MDA}{Architecture dirigée par les modèles (\textit{Model Driven Architecture})}

%\newabbreviation
%[type=\acronymtype]
%{bpm}{BPM}{\textit{Business Process Management}}

\newabbreviation
[type=\acronymtype]
{cmof}{CMOF}{Mise en place complète de méta-objets (\textit{Complete Meta-Object Facility})}

\newabbreviation
[type=\acronymtype]
{mof}{MOF}{Mise en place de méta-objets (\textit{Meta-Object Facility})}

%nom commun
\newabbreviation
[type=\acronymtype]
{icase}{I-CASE}{Environnement intégré de génie logiciel assisté par ordinateur (\textit{Integrated CASE Environment})}

%nom commun
\newabbreviation
[type=\acronymtype]
{alm}{ALM}{Gestion du cycle de vie des applications (\textit{Application lifecycle management})}

%nom propre
%\newabbreviation
%[type=\acronymtype]
%{acm}{ACM}{Association for Computing Machinery}

%\newabbreviation
%[type=\acronymtype]
%{usfda}{US FDA}{\textit{U.S. Food and Drug Administration}}

%\newabbreviation
%[type=\acronymtype]
%{dod}{DOD}{\textit{Departement Of Defense}}

%\newabbreviation
%[type=\acronymtype]
%{ieee}{IEEE}{\textit{Institute of Electrical and Electronics Engineers}}

\newabbreviation
[type=\acronymtype]
{emf}{EMF}{Cadre de modélisation Eclipse (\textit{Eclipse Modeling Framework})}

\newabbreviation
[type=\acronymtype]
{ocl}{OCL}{Langage de contraintes pour objets (\textit{Object Constraint Language})}

%nom commun
\newabbreviation
[type=\acronymtype]
{api}{API}{Interface de programmation d'application (\textit{Application Programme Interface})}

\newabbreviation
[type=\acronymtype]
{mbse}{MBSE}{Ingénierie système basée sur les modèles (\textit{Model Based Systems Engineering})}

\newabbreviation
[type=\acronymtype]
{slebok}{SLEBOK}{Corpus de connaissances en génie des langages logiciels (\textit{Software Language Engineering Body of Knowledge})}

%\newabbreviation
%[type=\acronymtype]
%{ansi}{ANSI}{\textit{American National Standards Institute}}

%\newabbreviation
%[type=\acronymtype]
%{ans}{ANS}{\textit{American Nuclear Society}}

%nom commun
\newabbreviation
[type=\acronymtype]
{adl}{ADL}{Langage de description d'architecture (\textit{Architecture Description Language})}

%\newabbreviation
%[type=\acronymtype]
%{cocomo}{COCOMO}{\textit{Constructive Cost Model}}



% La commande \glsaddallunused qui est juste avant \end{document} dans le fichier .tex principal fait en sorte que toutes les entrées 
% d'abréviations définies dans le fichier prelim sont incluses dans la liste
% si ce comportement n'est pas souhaité, mettre en commentaires.

% Si cette ligne est en commentaire, les abréviations seront listées dans la table des abréviations 
% si elles sont utilisées dans le texte seulement. Pour ce faire, vous devez utiliser la

% commande \gls(etiquette) à chaque fois dans le texte.
% en début de chapitre, utiliser \glsentryfull{edp} à la première utilisation.
% consulter glossariesbegin.pdf dans la documentation du paquetage glossaries 
% ou la page LaTeX/Glossary dans le wikibook sur Latex (https://en.wikibooks.org/wiki/LaTeX/Glossary)

% Précisions
% Don’t use \gls in chapter or section headings as it can have some unpleasant
% side-effects. Instead use \glsentrytext for regular entries and one of
% \glsentryshort, \glsentrylong or \glsentryfull for acronyms.
% Alternatively use glossaries-extra which provides special commands for use in
% section headings and captions, such as \glsfmtshort{⟨label⟩}.


\newabbreviation
[type=glossaire]
{llm_1}{Grand modèle de langage (GML)} {«~Modèle de langage constitué par apprentissage profond à partir de mégadonnées.~» (tirée de \cite{Gdt_llm_0})}

\newabbreviation
[type=glossaire]
{entiteinfo}{Entité informatique}{Système, sous-système, composant logiciel ou composant matériel en lien avec l’informatique. (inspirée de~: 
\cite{iso_24765_2017, Gdt_entite_0, Gdt_informatique_0})}

\newabbreviation
[type=glossaire]
{logiciel}{Logiciel}{Ensemble de suites d’instructions interprétables par un ordinateur. (inspirée de~: \cite{iso_24765_2017, Gdt_logiciel_0, Gdt_programme_0})}

\newabbreviation
[type=glossaire]
{composantlogiciel}{Composant logiciel}{Logiciel ou partie d’un logiciel. (inspirée de~: 
\cite{iso_24765_2017, Larousse_composant_0, Gdt_composantlogiciel_0})}

\newabbreviation
[type=glossaire]
{document}{Document}{Écrit servant d’information. (inspirée de~: \cite{iso_24765_2017, Larousse_document_0})}

\newabbreviation
[type=glossaire]
{documenteninformatique}{Documentation informatique}{Ensemble de documents associé à un ensemble d’entités informatiques. (inspirée de~: 
\ccite{clements_documenting_2014}[p.~12][iso_24765_2017][][Gdt_documentation_0][])}

\newabbreviation
[type=glossaire]
{procede}{Procédé}{Ensemble structuré d'activités visant un but logiquement déterminé par des antécédents et des conséquents où des intrants donnent lieu à des extrants. (inspirée de~: \cite{Bdl_procede_0, Gdt_processus_1, Gdt_processus_0, Robert_methode_0})}

\newabbreviation
[type=glossaire]
{procedededev}{Procédé de développement}{Procédé dont le but est de créer des entités informatiques. (inspirée de~: \cite{iso_24765_2017, Larousse_developpement_0, Gdt_sdlc_0,  pmi_guidefr_2017})}

\newabbreviation
[type=glossaire]
{processus}{Processus}{Instance d'un procédé. (inspirée de~: \cite{Gdt_cycle_0, Gdt_processus_1, Gdt_processus_0, pmi_guidefr_2017})}

\newabbreviation
[type=glossaire]
{artefact}{Artefact}{Intrant tangible ou extrant tangible à un processus. (inspirée de~:  \ccite{iso_19506_2012}[p.8][Gdt_artefact_0][])}

\newabbreviation
[type=glossaire]
{phasep}{Phase d'un projet}{Regroupement d’activités logiquement interreliées découlant de la gestion d’un projet. (inspirée de~: \cite{iso_24765_2017, Larousse_phase_0})}

\newabbreviation
[type=glossaire]
{phased}{Phase d'un procédé de développement}{Regroupement d’activités logiquement interreliées découlant d’un procédé de développement. (inspirée de~: \cite{iso_24765_2017, Larousse_phase_0})}

\newabbreviation
[type=glossaire]
{realiser}{Réaliser}{Rendre tangible ce qui est abstrait. (inspirée de~: \cite{Larousse_concret_0, Larousse_tangible_0, Robert_realiser_0})}


\newabbreviation
[type=glossaire]
{dettetechnique}{Dette technique}{Passif technique résultant de choix opportunistes dont le coût de mise à niveau augmente avec le temps. (inspirée de~: \cite{avgeriou_reducing_2016, holvitie_technical_2018, Larousse_passif_0})}

\newabbreviation
[type=glossaire,sort={etablir}]
{etablir}{Établir}{Créer les conditions pour soutenir la pérennité. (inspirée de~: \cite{Robert_etablir_0, Larousse_etablir_0})}

\newabbreviation
[type=glossaire]
{instrument}{Instrument}{«~Objet fabriqué servant à exécuter [quelque chose], à faire une opération.~» (tirée de \cite{Robert_instrument_0})}

\newabbreviation
[type=glossaire]
{methode}{Méthode}{«~Ensemble de démarches raisonnées, suivies pour parvenir à un but.~» (tirée de \cite{Robert_methode_0})}

\newabbreviation
[type=glossaire]
{technique}{Technique}{«~Ensemble des applications de la science dans le domaine de la production.~» (tirée de \cite{Larousse_technique_0})}

\newabbreviation
[type=glossaire]
{organiser}{Organiser}{«~Doter d’une structure, d’une constitution déterminée, d’un mode de fonctionnement.~» (tirée de \cite{Robert_organiser_0})}

\newabbreviation
[type=glossaire]
{module}{Module}{Est un composant comprenant une interface. (inspirée de~: \ccite{clements_documenting_2014}[p.~29-31][iso_24765_2017][p.~281])}

%\newabbreviation
%[type=glossaire]
%{cohesion}{Cohésion}{Mesure les liaisons selon une logique entre les éléments d'un même composant logiciel. (inspirée de~: \ccite{iso_26514_2021}[][lano_arch_2023][p.~24,237][Larousse_cohesion_0][])}

\newabbreviation
[type=glossaire]
{couplage}{Couplage}{Mesure l’interdépendance entre des composants logiciels. (inspirée de~:  \ccite{iso_26514_2021}[][lano_arch_2023][p.~237][Larousse_coupler_0][])}

\newabbreviation
[type=glossaire]
{adaptabilite}{Adaptabilité}{La capacité d'une entité à être modifiée afin de répondre à des changements d'environnements ou de besoins. (inspirée de~:  \ccite{bass_software_2021}[p.~117][Antidote_0][][iso_9126part1_2001][p.~10][Larousse_evolution_0][])}

\newabbreviation
[type=glossaire]
{langage}{Langage}{Système pour communiquer comprenant des symboles et des règles. (inspirée de~: \cite{Gdt_langage_0, Robert_langage_0, Robert_langue_0})}

\newabbreviation
[type=glossaire]
{langageformel}{Langage formel}{Est composé d’ensembles explicites de symboles et d’un ensemble explicite de règles appelé grammaire formelle. (inspirée de~: \ccite{hunter_metalogic_1973}[p.~4][iso_24765_2017][][Gdt_lf_0][][oregan_mathematics_2020][p.~188][roggenbach_formal_2021])}%[p.~6]

\newabbreviation
[type=glossaire]
{langagenaturel}{Langage naturel}{Est un langage non formel. (inspirée de~: \cite{iso_24765_2017, Gdt_ln_0})}

\newabbreviation
[type=glossaire]
{architecturelog}{Architecture informatique}{Organisation d'entités informatiques résultant de la phase de conception. (inspirée de~: \ccite{bass_software_2021}[p.~1,3][iso_42010_2022][p.~2][meyer_agile_2014][p.~37][printz_architecture_2012])}%[p.~61]

\newabbreviation
[type=glossaire]
{architecturedescription}{Langage de description d'architecture}{Langage avec une syntaxe et une sémantique spécifiées qui décrit l’architecture d’entité informatique. (inspirée de~: \ccite{iso_42010_2022}[p.~2,20])}

\newabbreviation
[type=glossaire]
{ontologie}{Ontologie informatique}{«~Corpus structuré de concepts, qui est modélisé dans un langage permettant l’exploitation par un ordinateur des relations sémantiques ou taxonomiques établies entre ces concepts.~» (tirée de \cite{Gdt_OntInfo_1})}

\newabbreviation
[type=glossaire]
{traitementautomatiquedes}{Traitement automatique des langues (TAL)}{«~Technique d’apprentissage automatique qui permet à l’ordinateur de comprendre le langage humain.~» (tirée de \cite{Gdt_NLP_0})}

\newabbreviation
[type=glossaire]
{apprentissagepro}{Apprentissage profond (AP)}{«~Mode d’apprentissage automatique généralement effectué par un réseau de neurones artificiels composé de plusieurs couches de neurones hiérarchisées.~» (tirée de \cite{Gdt_DL_0})}

\newabbreviation
[type=glossaire]
{boitenoire}{Boîte noire}{«~Système fermé dont on ignore le fonctionnement interne, qui n’est connu que par son comportement extérieur en fonction des sollicitations qui lui sont appliquées.~» \cite{Gdt_noire_0}}

\newabbreviation
[type=glossaire]
{ellipse}{Ellipse}{«~Omission d’un ou plusieurs mots dans une phrase qui reste cependant compréhensible.~» \cite{Robert_ellipse_0}}

\newabbreviation
[type=glossaire]
{modele}{Modèle}{Représentation d’un sujet pour en permettre le traitement. (inspirée de~: \ccite{caplat_model_2008}[p.28][Robert_modele_0][][gigch_system_1991][p.135][velten_mathematical_2024][p.3])}

\newabbreviation
[type=glossaire]
{systemelogique}{Système logique}{Est composé d’un langage formel et d’un ensemble de règles (axiome, proposition, logique) permettant d’inférer. (inspirée de~: \ccite{gabbay_logical_1994}[p.217][hunter_metalogic_1973][p.7][oregan_mathematics_2020])}%[p.211]

\newabbreviation
[type=glossaire]
{machineabstra}{Machine abstraite}{Système formel comprenant les instructions nécessaires au traitement automatique de données. (inspirée de~: \cite{Wolfram_abstractmachine_0,Larousse_ordinateur_0})}

\newabbreviation
[type=glossaire]
{modeleformel}{Modèle formel}{Est composé d’un langage formel et d’un ensemble de règles et vise la représentation d’un sujet pour en permettre le traitement. (inspirée de~: \ccite{abrial_modeling_2010}[p.176][abrial_refinement_2007][][davidsson_simulation_2017][p.785])}

\newabbreviation
[type=glossaire]
{coherent}{Cohérent}{Système logique où aucune contradiction n'a encore été découverte. (inspirée de~: \ccite{Gdt_coherence_0}[][Gdt_consistance_0][][oregan_mathematics_2020][p.214])}

\newabbreviation
[type=glossaire]
{couverture}{Couverture}{Ratio de la représentation d'une chose par une autre. Cela peut être d'un modèle par un artefact, d'un sujet par un modèle ou d'une théorie par un modèle. (inspirée de~: \cite{Antidote_0, Larousse_adequat_0, Larousse_adequation_0, Larousse_complet_0, Robert_exhaustif_0})}

\newabbreviation
[type=glossaire]
{completude}{Complétude}{Un système formel est complet lorsque toutes les vérités qu'il contient peuvent être déduites des axiomes et des règles d'inférence. (inspirée de~: \ccite{Larousse_completude_0}[][oregan_mathematics_2020][p.214])}

\newabbreviation
[type=glossaire]
{comprehensible}{Compréhensible}{Qui peut être intelligible par le lectorat. (inspirée de~: \ccite{krogstie_quality_2016}[p.58-60][Larousse_comprehensible_0][][olive_conceptual_2007][p.28-30])}

\newabbreviation
[type=glossaire]
{arithmetiquedepeano}{Arithmétique de Peano}{Système formel de l'arithmétique élémentaire qui permet des raisonnements par récurrence. (inspirée de~: \ccite{delahaye_godel_2025}[p.14,66,199])}

\newabbreviation
[type=glossaire]
{outildegenie}{Outil de génie logiciel assisté par ordinateur (\textit{CASE tool}, en anglais)}{Instrument servant à réaliser ou à établir des entités. (inspirée de~: \cite{iso_24765_2017,spurr_case_1990})}

\newabbreviation
[type=glossaire]
{meta}{Méta (\textit{meta} en grec)}{«~L’élément méta‑ vient du grec meta, qui peut signifier “au milieu”, “avec” ou “après”. \elide il exprime la succession, le changement, la participation, ou un concept qui englobe un autre concept.~» (tirée de \cite{Gdtmeta_1})}

\newabbreviation
[type=glossaire]
{langagespecifique}{Langage spécifique à un domaine (\textit{Domain-specific language}, en anglais)}{Langage destiné aux utilisateurs dans un domaine spécifique dont la syntaxe ou la sémantique peut être défini formellement. (inspirée de : \ccite{Gdt_dsl_0}[][wasowski_dsl_2023][p.26-36])}

\newabbreviation
[type=glossaire]
{automate}{Automate}{Cas particulier de machine abstraite formellement défini et conçu pour en permettre l'étude. (inspirée de~: \ccite{hopcroft_introduction_2007}[p.1][kandar_introduction_2013][p.48][Gdt_automate_0][])}

\newabbreviation
[type=glossaire]
{connaissance}{Connaissance}{Une croyance justifiée qui à force de légitimités devient un fait. (inspirée de~: \cite{gingras_qu_2026})}

\newabbreviation
[type=glossaire]
{savoireninfor}{Savoir en informatique appliquée}{Ensemble des connaissances s’appliquant à l’information, aux traitements de l’information et aux entités informatiques utilisé en pratique. (inspirée de~: \cite{Gdt_1_savoir_1, Gdt_1_savoir_2})}

\newabbreviation
[type=glossaire]
{gestiondelaconnaissance}{Gestion de la connaissance}{Organiser, contrôler et rendre accessibles les représentations des connaissances au bénéfice d’une organisation. (inspirée de~: \ccite{Gdt_km_0}[][pmi_guidefr_2017][p.136])}

\newabbreviation
[type=glossaire]
{informaticien}{Informaticien}{Spécialiste du traitement automatique et rationnel de l’information. (inspirée de~: \cite{Larousse_informatique_0,Gdt_informatique_0})}

\newabbreviation
[type=glossaire]
{informaticienapp}{Informaticien appliqué}{Informaticien spécialiste dans l'information provenant de domaines autres que l'informatique. (inspiré de : \cite{Gdt_inginf_0, Gdt_swe_0, Larousse_applique_0})}

\newabbreviation
[type=glossaire]
{praticieninformatique}{Praticien informatique}{Personne qui réalise ou établit des entités informatiques. (inspirée de~: \cite{Larousse_praticien_0, Gdt_praticien_0, Gdt_informatique_0})}

\newabbreviation
[type=glossaire]
{professionnelinformatique}{Professionnel informatique}{Personne exerçant une activité rémunérée liée à l'informatique. (inspirée de~: \cite{Larousse_professionnel_0, Gdt_informatique_0})}

\newabbreviation
[type=glossaire]
{participantinformatique}{Participant informatique}{Personne participant à des processus utilisant l’informatique. (inspirée de~: \cite{Larousse_participant_0, Larousse_participer_0, Gdt_informatique_0})}

\newabbreviation
[type=glossaire]
{taxinomie}{Taxinomie}{Activité et résultat de l’étude d’un sujet pour en obtenir des règles de classification. (inspirée de~: \cite{gdt_taxinomie_0, ralph-process-2019})}

\newabbreviation
[type=glossaire]
{retroingenierie}{Rétro-ingénierie}{Consiste à étudier un groupe d’artefacts pour en déterminer le fonctionnement et la méthode de fabrication. (inspirée de~: \cite{Gdt_RetroIng_1, iso_24765_2017})}

\newabbreviation
[type=glossaire]
{norme}{Norme}{«~Document, établi par consensus et approuvé par un organisme reconnu, qui fournit, pour des usages communs et répétés, des règles \elide~» (tirée de \ccite{iso_guide2_2004}[p.12])}

\newabbreviation
[type=glossaire]
{organismedeno}{Organisme de normalisation}{«~\elide principales fonctions, en vertu de ses statuts, est la préparation, l'approbation ou l'adoption de normes \elide.~» (tirée de \ccite{iso_guide2_2004}[p.12])}

\newabbreviation
[type=glossaire]
{reglement}{Règlement}{«~Document qui contient des règles à caractère obligatoire et qui a été adopté par une autorité.~» (tirée de \ccite{iso_guide2_2004}[p.16])}

\newabbreviation
[type=glossaire]
{standard}{Standard}{Document qui contient des règles de pratiques communes ne provenant pas d’un organisme de normalisation. (inspirée de~: \cite{iso_guide2_2004, Larousse_standard_0, Gdt_standard_0})}

\newabbreviation
[type=glossaire]
{convention}{Convention}{«~Accord de volonté entre deux ou plusieurs personnes physiques ou morales, par lequel elles s'engagent à faire ou à ne pas faire quelque chose.~» (tirée de \cite{Gdt_convention_0}}

\newabbreviation
[type=glossaire]
{patrimonial}{Patrimonial}{«~Se dit de composants, logiciels ou matériels, issus d'une génération ou d'une version dépassée et qui continuent d'être utilisés, après des ajustements, en même temps que la technologie contemporaine d'une entreprise.~» (tirée de \cite{Gdt_patrimonial_0})}






